\documentclass[11pt,a4paper]{article}
\usepackage{../../intuitmath}
\hypersetup{pdftitle={Why Hamiltonian Mechanics?},pdfauthor={Mathematics and Physics: Intuition First}}
\renewcommand{\sectionmark}[1]{\markboth{#1}{}}

\begin{document}
\tutorialtitle{Why Hamiltonian Mechanics?}{Predicting the next instant from position and momentum}{Discovering Hamilton's equations before deriving them}
\tableofcontents
\clearpage

\section{The question that motivates the formalism}

Lagrangian mechanics describes a configuration using $q$ and a velocity $\dot q$. It also introduces the conjugate momentum
\[
p=\pd{L}{\dot q}.
\]
For an ordinary particle, $p=m\dot q$, so position and momentum contain the same instantaneous information as position and velocity.

This suggests a practical question:

\begin{intuition}[What we actually need]
If the particle is currently at $(q,p)$, how do we find its position and momentum a moment later?
\end{intuition}

For a short interval $\Delta t$,
\[
q(t+\Delta t)\approx q(t)+\dot q\,\Delta t,
\qquad
p(t+\Delta t)\approx p(t)+\dot p\,\Delta t.
\]
We therefore need two rates:
\[
\dot q=f(q,p,t),
\qquad
\dot p=g(q,p,t).
\]
No partial derivatives of the Hamiltonian have been assumed. We have merely identified the information required to evolve the state.

\section{Discovering the equations for a familiar particle}

For a particle in a potential $V(q)$, define
\[
H(q,p)=\frac{p^2}{2m}+V(q).
\]
We already know
\[
\dot q=\frac pm,
\qquad
\dot p=F=-\td{V}{q}.
\]
Now differentiate $H$:
\[
\pd{H}{p}=\frac pm=\dot q,
\qquad
\pd{H}{q}=\td{V}{q}=-\dot p.
\]

\begin{resultbox}[Hamilton's equations, discovered]
\[
\boxed{
\dot q=\pd{H}{p},
\qquad
\dot p=-\pd{H}{q}.}
\]
The two rates needed to predict the next instant are already contained in the two partial derivatives of a single function.
\end{resultbox}

\section{Why this is significant}

Newton's equation is usually second order:
\[
m\ddot q=F(q).
\]
Hamilton's formulation replaces it with two coupled first-order equations for the state $(q,p)$. This is not less information: specifying $q$ and $p$ is equivalent to specifying $q$ and $\dot q$ for a regular mechanical system.

The pair $(q,p)$ is a point in \emph{phase space}. As time changes, the state traces a curve there.

\begin{center}
\begin{tikzpicture}[x=1.25cm,y=1.0cm,>=Latex]
  \draw[->,ink!65] (-3.1,0)--(3.35,0) node[right]{$q$};
  \draw[->,ink!65] (0,-2.4)--(0,2.65) node[above]{$p$};
  \draw[blue,line width=1.6pt,->] (2.55,0) arc[start angle=0,end angle=325,x radius=2.55,y radius=1.75];
  \fill[teal] (1.35,1.49) circle (3pt) node[above right]{$(q,p)$};
  \draw[->,gold,line width=1.2pt] (1.35,1.49)--++(-.75,.32) node[above left]{$(\dot q,\dot p)$};
  \node[ink] at (2.35,2.15){one oscillator orbit};
\end{tikzpicture}
\end{center}

At every point, Hamilton's equations assign the tangent vector
\[
(\dot q,\dot p)=\left(\pd{H}{p},-\pd{H}{q}\right).
\]
The Hamiltonian is therefore a rule for the flow through phase space.

\section{Why change variables from velocity to momentum?}

The Lagrangian is naturally a function of $(q,\dot q,t)$. We want a new function of $(q,p,t)$. Merely replacing symbols is not enough because $p$ is defined by a derivative of $L$.

The required change is the Legendre transform:
\[
\boxed{H(q,p,t)=p\dot q-L(q,\dot q,t),}
\]
where $\dot q$ is re-expressed in terms of $(q,p,t)$ using
\[
p=\pd{L}{\dot q}.
\]

\begin{intuition}[What the Legendre transform does]
It replaces the slope-like variable $\dot q$ by the conjugate slope $p$. The correction term $p\dot q$ preserves the information while changing which variable is independent.
\end{intuition}

For
\[
L=\frac12m\dot q^2-V(q),
\qquad p=m\dot q,
\]
we have $\dot q=p/m$ and
\begin{align*}
H&=p\frac pm-\left[\frac12m\left(\frac pm\right)^2-V(q)\right]\\
&=\frac{p^2}{2m}+V(q)=T+V.
\end{align*}
This familiar equality relies on ordinary quadratic kinetic energy and a velocity-independent potential; $H$ is not automatically ``the energy'' in every model.

\section{General derivation of Hamilton's equations}

Begin with
\[
H=p\dot q-L.
\]
Take a total differential:
\[
\dd H=\dot q\,\dd p+p\,\dd\dot q-\dd L.
\]
Since
\[
\dd L=\pd{L}{q}\dd q+\pd{L}{\dot q}\dd\dot q+\pd{L}{t}\dd t
\]
and $p=\partial L/\partial\dot q$, the two terms containing $\dd\dot q$ cancel:
\[
\dd H=\dot q\,\dd p-\pd{L}{q}\dd q-\pd{L}{t}\dd t.
\]
The Euler--Lagrange equation gives
\[
\pd{L}{q}=\dot p.
\]
Therefore
\[
\dd H=\dot q\,\dd p-\dot p\,\dd q-\pd{L}{t}\dd t.
\]
But $H=H(q,p,t)$ also implies
\[
\dd H=\pd{H}{q}\dd q+\pd{H}{p}\dd p+\pd{H}{t}\dd t.
\]
Comparing independent coefficients yields
\[
\boxed{
\dot q=\pd{H}{p},
\qquad
\dot p=-\pd{H}{q},
\qquad
\pd{H}{t}=-\pd{L}{t}.}
\]

\section{Many degrees of freedom}

For coordinates $q_i$ and momenta
\[
p_i=\pd{L}{\dot q_i},
\]
define
\[
H(\vect q,\vect p,t)=\sum_i p_i\dot q_i-L.
\]
Then
\[
\boxed{
\dot q_i=\pd{H}{p_i},
\qquad
\dot p_i=-\pd{H}{q_i}.}
\]
The phase space has dimension $2n$: one coordinate and one conjugate momentum for each degree of freedom.

\section{Worked example: harmonic oscillator}

For
\[
H(q,p)=\frac{p^2}{2m}+\frac12kq^2,
\]
Hamilton's equations give
\[
\dot q=\frac pm,
\qquad
\dot p=-kq.
\]
Differentiate the first equation and insert the second:
\[
\ddot q=\frac{\dot p}{m}=-\frac kmq.
\]
Thus
\[
\ddot q+\omega^2q=0,
\qquad \omega=\sqrt{k/m}.
\]

The constant-energy curves satisfy
\[
\frac{p^2}{2m}+\frac12kq^2=E,
\]
which are ellipses in the $(q,p)$ plane. The oscillator moves around one ellipse without crossing to another energy.

\section{Conservation of the Hamiltonian}

Along a trajectory,
\begin{align*}
\td{H}{t}
&=\pd{H}{q}\dot q+\pd{H}{p}\dot p+\pd{H}{t}\\
&=(-\dot p)\dot q+(\dot q)\dot p+\pd{H}{t}\\
&=\pd{H}{t}.
\end{align*}
Therefore, if $H$ has no explicit time dependence,
\[
\boxed{H=\text{constant along the motion}.}
\]

\section{Poisson brackets}

For functions $A(q,p,t)$ and $B(q,p,t)$ define
\[
\{A,B\}=\pd{A}{q}\pd{B}{p}-\pd{A}{p}\pd{B}{q}.
\]
Hamilton's equations can be written
\[
\dot q=\{q,H\},
\qquad
\dot p=\{p,H\}.
\]
More generally,
\[
\boxed{\td{A}{t}=\{A,H\}+\pd{A}{t}.}
\]
If $A$ has no explicit time dependence and $\{A,H\}=0$, then $A$ is conserved.

The basic bracket is
\[
\{q,p\}=1.
\]
This compact algebra becomes the classical precursor of commutators in quantum mechanics.

\section{Phase-space geometry}

In two-dimensional phase space the velocity field is
\[
\vect v_{\rm ph}=\left(\pd{H}{p},-\pd{H}{q}\right).
\]
Its divergence is
\[
\pd{}{q}\left(\pd{H}{p}\right)
+\pd{}{p}\left(-\pd{H}{q}\right)=0
\]
when mixed derivatives commute. Small phase-space areas are neither compressed nor expanded by Hamiltonian flow. This is the local heart of Liouville's theorem.

\section{When the Legendre transform can fail}

To solve $p=\partial L/\partial\dot q$ for $\dot q$, the velocity Hessian must be invertible. In one dimension this means
\[
\frac{\partial^2L}{\partial\dot q^2}\neq0.
\]
If it is singular, the system contains constraints and needs a more careful Hamiltonian treatment. The standard equations remain correct for regular Lagrangians.

\section{A numerical viewpoint}

Hamilton's equations are an initial-value system. Given $(q_0,p_0)$, a simple first step is
\[
q_{n+1}=q_n+\Delta t\,\pd{H}{p}(q_n,p_n),
\qquad
p_{n+1}=p_n-\Delta t\,\pd{H}{q}(q_n,p_n).
\]
This explicit Euler rule is easy but does not preserve Hamiltonian geometry well over long times. Symplectic methods alter the update order to respect phase-space structure more faithfully.

\section{What to remember}

\begin{enumerate}
\item Position and momentum specify the instantaneous state.
\item Predicting the next instant requires $\dot q$ and $\dot p$.
\item For a familiar particle, these rates are the two partial derivatives of $H$ with opposite signs.
\item The Legendre transform proves that the same structure follows generally from a regular Lagrangian.
\item The Hamiltonian generates a flow through phase space.
\end{enumerate}

\begin{checkBox}[The central insight]
Hamilton's equations are not arbitrary formulas involving partial derivatives. They say that a single scalar function contains the two local rates needed to move the complete state $(q,p)$ forward in time.
\end{checkBox}

\end{document}
